% Options for packages loaded elsewhere
\PassOptionsToPackage{unicode}{hyperref}
\PassOptionsToPackage{hyphens}{url}
%
\documentclass[
  11pt,
]{article}
\usepackage{amsmath,amssymb}
\usepackage{iftex}
\ifPDFTeX
  \usepackage[T1]{fontenc}
  \usepackage[utf8]{inputenc}
  \usepackage{textcomp} % provide euro and other symbols
\else % if luatex or xetex
  \usepackage{unicode-math} % this also loads fontspec
  \defaultfontfeatures{Scale=MatchLowercase}
  \defaultfontfeatures[\rmfamily]{Ligatures=TeX,Scale=1}
\fi
\usepackage{lmodern}
\ifPDFTeX\else
  % xetex/luatex font selection
\fi
% Use upquote if available, for straight quotes in verbatim environments
\IfFileExists{upquote.sty}{\usepackage{upquote}}{}
\IfFileExists{microtype.sty}{% use microtype if available
  \usepackage[]{microtype}
  \UseMicrotypeSet[protrusion]{basicmath} % disable protrusion for tt fonts
}{}
\makeatletter
\@ifundefined{KOMAClassName}{% if non-KOMA class
  \IfFileExists{parskip.sty}{%
    \usepackage{parskip}
  }{% else
    \setlength{\parindent}{0pt}
    \setlength{\parskip}{6pt plus 2pt minus 1pt}}
}{% if KOMA class
  \KOMAoptions{parskip=half}}
\makeatother
\usepackage{xcolor}
\usepackage[margin=25mm]{geometry}
\setlength{\emergencystretch}{3em} % prevent overfull lines
\providecommand{\tightlist}{%
  \setlength{\itemsep}{0pt}\setlength{\parskip}{0pt}}
\setcounter{secnumdepth}{-\maxdimen} % remove section numbering
\ifLuaTeX
  \usepackage{selnolig}  % disable illegal ligatures
\fi
\usepackage{bookmark}
\IfFileExists{xurl.sty}{\usepackage{xurl}}{} % add URL line breaks if available
\urlstyle{same}
\hypersetup{
  pdftitle={A general reference moment and intrinsic finite-energy targets},
  pdfauthor={Unverified research draft, relative to seven complete caller grants},
  hidelinks,
  pdfcreator={LaTeX via pandoc}}

\title{A general reference moment and intrinsic finite-energy targets}
\usepackage{etoolbox}
\makeatletter
\providecommand{\subtitle}[1]{% add subtitle to \maketitle
  \apptocmd{\@title}{\par {\large #1 \par}}{}{}
}
\makeatother
\subtitle{New partial continuation of original P16.4; the
prescribed-target equation remains open}
\author{Unverified research draft, relative to seven complete caller
grants}
\date{7 October 2026}

\begin{document}
\maketitle

\subsection{Status, scope, and the material
correction}\label{status-scope-and-the-material-correction}

This paper does not close original P16.4. Its new result is a proposed
general-dimensional uniform reference moment and the resulting intrinsic
target criterion. The proof below reconstructs the geometric
applications; ROOT248, ROOT251, ROOT252, native243, and whole250 are not
grants. Successful finite checks and compilation do not verify the
mathematics. The source ledger specifies the ordinary foundational
results used and the boundary of direct source reading.

Retain a complete real-valued non-Archimedean field \(K\) with
characteristic and residue characteristic zero, a smooth geometrically
integral projective \(K\)-variety \(X\) of dimension \(n\geq1\), and an
actual ample rational line bundle \(L\). Valuations may be trivial,
discrete, or nondiscrete, and \(K\) need not be algebraically closed.
Set \(T=X^{\mathrm{an}}\) and \(V=(L^n)>0\). All measures have
\textbf{raw intersection mass} \(V\). All functions are evaluated on the
entire analytic carrier, including points of nongeneric scheme support
and literal poles. No measure-support or invariance assumption is made.
Singular varieties, positive residue characteristic, nonample classes,
nonprojectivity, and nonalgebraizability remain separate axes.

Exactly O217, D214, N220, T171, Q222, C218, and G227, with their
unchanged clarifications, are supplied mathematical grants. O supplies
the actual finite-energy mixed operator and energy calculus. D supplies
completed-algebraic-closure descent and the invariant Radon lift. T
supplies the full trivial-valuation equation. G227 supplies the full
original-field genus-zero equation, including nonsplit conics; that
theorem is used as a whole and is not reproved. N, Q, and C retain their
stated special hypotheses. C gives locality and comparison, not an
inverse theorem. The remaining arguments here are new drafts using
explicitly stated ordinary source theorems.

Whole250 failed independent review198. Its Theorem8 and note n4 asserted
coercivity without the necessary hypothesis \(U_\mu<\infty\). That
assertion is false. Adding the hypothesis is a \textbf{material change
in target scope}, not a harmless correction. The failed paper, argument,
ledger, native final, and failure report remain unchanged as input data.
No other whole250 claim is promoted to a grant by the reviewer's
discussion.

\textbf{Proposition 1 (actual counterexample to the failed scope).} On
\(\mathbb P^1_C\) with \(L=\mathcal O(1)\) and the standard Gauss
reference, the uniform reference moment is zero, but an arbitrary target
of correct raw mass need not yield coercivity.

\textbf{Proof.} Take complete algebraically closed nontrivial \(C\) in
the retained characteristic domain. Let \(\zeta\) be its Gauss point and
\(p=(z=0)\). Choose \(0<|a|<1\), and put \(\ell_m=-m\log|a|\). The
generating sections \(T_1,a^mT_0\) give the actual continuous FS scalar
\[
 u_m=\log\max(|z|,e^{-\ell_m})-\log\max(1,|z|).
\] It has supremum zero. N gives \(\operatorname{MA}(0)=\delta_\zeta\)
and \(u(\zeta)=\sup u\) for every actual psh scalar. Thus \(M_0=0\). Q
identifies the genuine curvature of \(u_m\) as \(\delta_{y_m}\), where
\(y_m\) is the disk Gauss point of radius \(|a|^m\). Both \(u_m(p)\) and
\(u_m(y_m)\) are \(-\ell_m\). O's degree-one energy identity gives
\(E(u_m)=-\ell_m/2\). For \(\mu=\delta_p\), \[
 B(u_m)=\ell_m/2\longrightarrow\infty,\qquad
 E(u_m)-\int u_m\,d\mu=\ell_m/2\longrightarrow+\infty.
\] O's real segment calculus gives \(B(tu_m)=t^2\ell_m/2\). Taking
\(t=\ell_m^{-1/2}\) for large \(m\) gives cost \(1/2\) and target
integral \(\sqrt{\ell_m}\), so \(U_\mu=\infty\). The unit bound is
therefore essential. These are genuine metrics and genuine MA measures;
no graph measure is substituted. \(\square\)

\subsection{Actual operator, energy, and the new
statement}\label{actual-operator-energy-and-the-new-statement}

Fix a continuous psh reference \(\theta\) on \(L\). An actual psh metric
means the all-point decreasing-FS-net cone of BE Definition7.2, with the
identically \(-\infty\) endpoint excluded. For continuous psh scalars,
put \[
 E_\theta(v)=\frac1{n+1}\sum_{j=0}^n\int_T v\,
 \operatorname{MA}(v^{[j]},0^{[n-j]}).
\] For an actual endpoint, define \(E_\theta(u)\) as the infimum of
energies of continuous actual psh majorants. This extended energy is
used even before nontriviality or finite-energy membership is
established. Define \(\mathcal E^1=\{u:E_\theta(u)>-\infty\}\). Grant O
supplies the positive genuine mixed operator on \((\mathcal E^1)^n\),
approximation independence, arbitrary decreasing-net continuity,
all-test integrability, the genuine energy identity, affine-segment
derivative, reference covariance and complete-field pushforward. Each
output has mass \(V\).

Write \(\sup_Tu=0\) for normalization. Then \(u\leq0\) and
\(B(u)=-E_\theta(u)\geq0\). For a positive Radon target of mass \(V\),
define \[
 A_\mu(u)=\int_T(-u)\,d\mu,\qquad
 U_\mu=\sup_{\substack{u\in\mathcal E^1,\ \sup u=0\\B(u)\leq1}}A_\mu(u),
\] initially allowing infinity, and \[
 J_\theta(\mu)=\sup_{v\text{ continuous actual psh}}
 \left(E_\theta(v)-\int_Tv\,d\mu\right).
\] Constants cancel exactly because the mass is \(V\).

\textbf{Theorem 2 (new reference moment and target criterion, draft
proof).} In the retained domain, one can choose an original-field
zero-weight FS reference \(\theta\) and a finite constant \(M_0\) such
that \[
 \int_T(-u)\,\operatorname{MA}_\theta(0)\leq M_0
 \quad\text{for every actual psh }u\text{ with }\sup u=0.
 \tag{2.1}
\] For every positive Radon \(\mu\) of raw mass \(V\), the following
conditions are equivalent: every actual \(\mathcal E^1\) potential is
absolutely \(\mu\)-integrable; \(U_\mu<\infty\); and
\(J_\theta(\mu)<\infty\). When these finite target conditions hold, with
\(D_n=n(n+1)/2\), \[
 A_\mu(u)\leq U_\mu\bigl(1+2M_0+\sqrt{2D_nB(u)}\bigr),
 \tag{2.2}
\] \[
 E_\theta(u)-\int_Tu\,d\mu
 \leq-B(u)+U_\mu(1+2M_0)+U_\mu\sqrt{2D_nB(u)},
 \tag{2.3}
\] \[
 J_\theta(\mu)\leq U_\mu(1+2M_0)+\tfrac12D_nU_\mu^2.
 \tag{2.4}
\] In particular, \textbf{for such finite-energy targets}, the objective
tends to \(-\infty\) along normalized costs tending to infinity. Its
supremum over actual \(\mathcal E^1\) equals its continuous-test
supremum. No prescribed-target attainment or general every-point
uniqueness is claimed.

For nontrivial \(K\), the construction below gives \(M_0=VaD_0\), where
\(a=d^{n-1}V\) and \(D_0\) is the diameter between finitely many
specified points on one fixed actual auxiliary-curve skeleton. For
\(n=1\), the original curve over the completed algebraic closure is used
and \(a=V\). For trivial \(K\), T supplies the criterion and equation;
the canonical reference has \(M_0=0\). The assertion for the selected
reference is sufficient: reference covariance makes \(J\) finiteness and
the absolute energy class intrinsic. We do not assert that every other
reference has the same constant or atomic support.

\subsection{Source foundations and precise
boundaries}\label{source-foundations-and-precise-boundaries}

The following are ordinary source theorems, not additional caller
grants. Their immediate printed units and applications are read and
recorded in the ledger. Their entire recursively cited foundations are
not freshly reproved here.

\begin{itemize}
\tightlist
\item
  For a nontrivial complete field, a flat finite-type scheme model is
  finitely presented (BE4.5). Over an algebraically closed such field,
  normalization of a reduced generic fiber in its model is finite, and
  the normalized special fiber is reduced (BE4.20). The completed
  algebraic closure is algebraically closed and has divisible value
  group; these standard valued-field facts also occur in the supplied
  descent setup.
\item
  A strict affinoid algebra has finite-type reduced canonical reduction;
  strict Noether normalization gives a finite injective Tate-polydisc
  algebra of the analytic dimension. Canonical reductions preserve
  equidimension, and their map to any admissible formal reduction is
  finite and surjective (GRW2.10,2.12). The canonical reduction map is
  surjective; each generic reduction point has exactly one inverse
  point, and these are precisely the finite Shilov boundary (GRW2.16,
  referring to Berkovich2.4.4). Laurent localization at a norm-one
  element reduces to localization of the canonical reduction. These
  statements hold without a Noetherian valuation-ring assumption.
\item
  The analytification of a smooth pure-dimensional variety has the same
  local dimension, and its affinoid domains are pure-dimensional
  (GRW2.4). At a Shilov point of a strict affinoid, residual
  transcendence degree equals local analytic dimension (Poineau4.4). The
  latter immediate proof isolates a canonical component by Laurent
  localization and uses its function-field embedding plus Abhyankar's
  upper bound.
\item
  Finite pullback of an ample invertible sheaf is ample: use the
  quasi-affine forward theorem Stacks29.38.7(2), Tag0892.
  Smooth-projective field Serre duality is
  \(H^q(A^{-r})^\vee=H^{n-q}(\omega_X\otimes A^r)\); it follows from
  Stacks48.27.1(5),(7), Tag0FVV, for invertible \(A\). Positive Serre
  vanishing is Stacks30.16.1(4), Tag0B5T. These are applied over fields,
  where the Noetherian hypotheses are satisfied.
\item
  A smooth projective geometrically integral curve over a complete
  algebraically closed nontrivial rank-one field admits a semistable
  vertex set (BPR4.22). Any finite type-II set can be included, and
  deleting finitely many algebraic points permits a refinement retaining
  every old edge with its old logarithmic length (BPR3.13). This uses
  the ordinary stable reduction theorem, whose analytic rank-one
  application is stated in BPR's proof via BL85 Theorem7.1. Universal
  formal realization of arbitrary polarized graph data is not a premise.
\item
  A nonconstant rational function on a smooth projective curve defines a
  finite map to \(\mathbb P^1\) of degree equal to its pole-divisor
  degree (Stacks53.2.1--4, Tag0BXX, and33.44.9,11, Tags0AYY,0AYZ). An
  annulus has the usual Gauss norm, residue \(\widetilde B(\tau)\) at
  value-group radii, and logarithmic path length (BPR2.1--3).
\end{itemize}

This paper relies on those stated foundational results in their exact
domains. BGR Tate Noetherianity, strict Noether normalization,
powerbounded integrality, FK2017 closed-ideal foundations, Berkovich
generic-reduction correspondence, analytic dimension, the recursive
Serre-duality construction, and classical stable reduction have not all
been independently reconstructed from their ultimate sources. This is a
source-verification boundary, not a hidden geometric antecedent or a
claimed new proof of those foundations. In particular, the ledger does
not claim full recursive source certification. The theorem and source
applications require independent review. Direct-formal GAGA recognition
from whole250 remains unresolved route data: review198 premises14/15 and
their comparison antecedents are not proved or used here.

\subsection{One unchanged reference and all initial
points}\label{one-unchanged-reference-and-all-initial-points}

Assume \(K\) is nontrivial and put \(C=\widehat{\overline K}\). Choose
an original integral very ample multiple \(M=dL\), with an original
embedding basis \(S_0,\ldots,S_N\). Increase \(d\) until \[
 H^q(X_C,\mathcal O(-kM))=0
 \quad(1\leq k,q\leq n-1). \tag{4.1}
\] Indeed start with an integral ample multiple \(A\) of \(L\). Field
Serre duality and positive vanishing for \(\omega_X\otimes A^r\) kill
\(H^q(A^{-r})\) for \(1\leq q<n\) and all sufficiently large \(r\).
Choose a large very ample \(A^r\); every \(kr\) in the finite list is
beyond this threshold. Flat field base change preserves the vanishings.
The cohomology statements concern the actual original variety and
persist to \(C\), not reductions modulo a valuation element. For \(n=1\)
the list is empty.

Freeze \(\theta=\operatorname{FS}(S_0,\ldots,S_N)/d\) with zero weights.
Let \(\mathcal X_0\) be the schematic closure of \(X_C\) in
\(\mathbb P^N_{C^\circ}\) under this same embedding. On each affine
chart, the closure ring injects into its generic ring, so it is
torsion-free and flat over the valuation ring. It is finite type, and
BE4.5 makes it finitely presented. BE4.20 makes the normalization
\(p:\mathcal X\to\mathcal X_0\) finite with reduced special fiber. Its
generic morphism is the identity of \(X_C\), since \(X_C\) is normal.
The normalized model is again flat and finite type by BE4.20, so BE4.5
gives finite presentation for it too. The pullback
\(\mathcal M=p^*\mathcal O(1)\) is ample by the forward finite-pullback
theorem. A finite morphism is projective, so this is a projective model
suitable for BE6.3.

The original sections still generate \(\mathcal M\). In a model frame
their coefficients and a unit-ideal relation are bounded by one on the
formal generic domain. The relation forces the coefficient maximum to be
at least one, while integrality forces it to be at most one. Thus their
maximum is exactly one at \textbf{every} analytic point. The model
metric divided by \(d\) is exactly \(\theta_C\); no vertical correction,
carrier replacement, or arbitrary lattice-norm extension identity is
used.

Here the formal completion really is topologically finitely presented,
as required later by CLD6.5.1. Fix \(0<|a|<1\) in \(C\) and write a flat
affine model chart as \(R[t_1,\ldots,t_m]/(f_1,\ldots,f_r)\),
\(R=C^\circ\). The ring \(R\) is \(a\)-adically complete. Every
restricted-power-series algebra \(R\langle t\rangle\) is Noetherian
outside \((a)\), because inverting \(a\) gives the Noetherian strict
Tate algebra \(C\{t\}\). Thus \(R\) is topologically universally
Noetherian outside \((a)\). FK2017 Proposition0.8.4.4 says that every
ideal of \(R\langle t\rangle\) is closed in this domain. In particular
the ideal generated by the finite \(f_j\) is closed, and its quotient is
complete. Modulo \(a^j\) restricted series are polynomials, so this
quotient has exactly the same reductions as the original affine chart
and is its completion. It is topologically finitely presented. BE4.3
supplies flatness of this completion and identifies its generic domain
and reduction with those of the model. This uses a specific ordinary
restricted-series closed-ideal theorem; it is separate from, and
supplies no GAGA or coherent-section comparison from, the unresolved
direct-formal route.

Discard formal affine charts with empty generic domain; reduction
surjectivity makes their special fibers empty too. Every remaining
strict generic domain is a nonempty affinoid domain of the smooth
\(X_C\) and is pure \(n\)-dimensional. GRW2.10(3) gives pure
\(n\)-dimensional canonical reduction. Its proof uses strict Noether
normalization for total dimension, then norm-one Laurent localization
isolating each component for purity. GRW2.12 gives finite surjective
canonical-to-formal reduction: a Tate presentation and powerbounded
integrality give integrality, and its polynomial-relation argument
proves the reduction kernel nil, giving surjectivity.

A finite surjection from a finite-type pure \(n\)-dimensional field
scheme has pure \(n\)-dimensional target. Each source component has
closed image of the same dimension because its function-field extension
over that image is finite. Their finitely many images cover the target;
every maximal irreducible subset of this union is one such image.
Conversely each such \(n\)-dimensional image is a target component. Thus
\textbf{every} special-fiber component of \(\mathcal X\) has dimension
\(n\). A point mapping to a generic formal reduction point maps to a
generic canonical reduction point: finiteness makes the fiber over the
generic point zero-dimensional, and purity makes each of those points
generic. GRW2.16 therefore identifies these points with the finite
Shilov set on that chart. Their union is a finite nonempty set
\(S=\{y_i\}\).

The reduction composition is \[
 \operatorname{red}_{\rm formal}=\iota\circ
 \operatorname{red}_{\rm canonical}.
\] GM A.3 prints the opposite composition; that is a direction typo and
is not propagated. No singleton-fiber claim is used to identify an
entire residue field.

Poineau4.4 gives
\(\operatorname{trdeg}_{\widetilde C}\widetilde{\mathcal H(y_i)}=n\).
Each \(y_i\) has \textbf{generic scheme support in the original
\(X_C\)}. To prove this, let \(\xi_i\) be its scheme support, so
\(\mathcal H(y_i)\) is the completion of the valued \(\kappa(\xi_i)\).
Approximate unit lifts of \(n\) algebraically independent residues by
elements of \(\kappa(\xi_i)\) within error less than one. Their residues
persist. A nonzero polynomial relation over \(C\), divided by a
coefficient of maximal absolute value, would reduce to a nonzero
polynomial relation among those residues. Thus the lifts are
algebraically independent over \(C\), and
\(\operatorname{trdeg}_C\kappa(\xi_i)\geq n\). On the integral
\(n\)-dimensional variety only the generic point has that dimension.
Hence \(\mathcal H(y_i)\) is the completion of honest \(C(X)\).

Choose one original integral basis change as follows. At each
special-fiber generic point at least one original section is nonzero.
The coefficient vectors in \(\widetilde K^{N+1}\) giving a vanishing
combination there form a proper \(\widetilde K\)-linear subspace, even
when the component field extends \(\widetilde K\). A finite union of
proper subspaces over an infinite field cannot fill the vector space:
contain each in a hyperplane, whose product is a nonzero polynomial, and
evaluate it off its zero set. Residue characteristic zero makes
\(\widetilde K\) infinite. Lift a vector outside the union to
\(K^\circ\). A unit entry permits an invertible integral basis
completion, with integral inverse. The two ultrametric inequalities for
the matrix and its inverse preserve the section maximum at every
analytic point. Rename its first section \(S_0\). The reference
\(\theta\) is unchanged, \(S_0\) has model norm one at every \(y_i\),
and \[
 g_l=S_l/S_0,\qquad |g_l(y_i)|\leq1. \tag{4.2}
\]

On the original \(S_0\)-chart, the special coordinate algebra of
\(\mathcal X_0\) is generated by the reductions of these original
\(g_l\). A component \(Y\) of \(\mathcal X_s\) maps finitely onto its
\(n\)-dimensional integral closed image \(Z\) in that chart. Therefore
\(\kappa(Y)\) is finite over \[
 \widetilde C(\bar g_1,\ldots,\bar g_N)=\kappa(Z),
\] and this \textbf{coordinate subfield} has transcendence degree \(n\).
Reduction embeds it in \(\widetilde{\mathcal H(y_i)}\). Equality with
the entire latter field is unnecessary and is not asserted.

BE6.3(i), applied to each power \(\mathcal M^k\), gives \[
 \sup_{X_C^{\rm an}}|t|_{kd\theta_C}
       =\max_i|t(y_i)|_{kd\theta_C}
 \quad(t\in H^0(X_C,kM)). \tag{4.3}
\] This is section norming for every later section. After denominator
alignment, finite maxima with arbitrary real constant weights commute
with these suprema. Thus every later FS scalar is normed by the same
\(S\). This alone supplies no individual lower bound.

\subsection{Genuine atomic reference support on the same
model}\label{genuine-atomic-reference-support-on-the-same-model}

The formal completion is flat, separated, and locally topologically
finitely presented, so CLDv1 Lemma6.5.1 applies. At any point \(x\),
write \(W=\overline{\operatorname{red}_{\mathcal X}(x)}\) with its
reduced structure. It is a proper integral variety because
\(\mathcal X_s\) is proper. The lemma describes the germ reduction by
this \(W\). CLD6.6.8 then identifies the residual metric bundle with the
pullback of \(\mathcal M|_W\), using its actual integral transition
cocycle. Consequently its defining dimension is at most \(\dim W\).

If \(x\) reduces to a nongeneric special-fiber point, then \(\dim W<n\),
since every residual component has dimension \(n\). CLD6.9.2 puts the
top-wedge support in the set of defining dimension \(n\). Its proof,
via6.8.3, kills the local top tropical wedge in lower defining
dimension. Thus the support is contained in the points reducing to
component generic points, namely the same finite \(S\). The domain is
good, separated, pure-dimensional, and boundaryless. The rational bundle
is handled by taking the integral \(M\) and dividing the top wedge by
\(d^n\). O agrees with the genuine continuous CLD operator and supplies
positivity and raw mass \(V\). Hence \[
 \operatorname{MA}_{\theta_C}(0)=\sum_i b_i\delta_{y_i},
 \qquad b_i\geq0,\quad\sum_i b_i=V. \tag{5.1}
\] No \(b_i\) is calculated. CLD6.9.2 is an atomic-support theorem; the
coefficient theorems6.9.3--4 are not substituted for it or required
here. This avoids the unresolved raw-component-coefficient route of
whole250.

\subsection{One actual coefficient-Gauss incidence
curve}\label{one-actual-coefficient-gauss-incidence-curve}

Suppose \(n\geq2\), and set \(s=n-1\). For each \(i\), take the
unit-radius Gauss completion \[
 H'_i=\widehat{\mathcal H(y_i)(A_{jl}:1\leq j\leq s,\ 1\leq l\leq N)}.
\] Define \(A_{j0}=-\sum_{l=1}^N A_{jl}g_l(y_i)\) in \(H'_i\). We prove
that the full \(s(N+1)\) coefficient residues are independent over
\(k=\widetilde C\).

The original residual coordinate field is the function field of the
integral \(n\)-dimensional affine image \(Z_i\). Its smooth locus is
nonempty in characteristic zero. At a smooth point its coordinate
differentials \(d\bar g_l\) span the cotangent space. For the algebraic
map \[
 (z,a_{\rm free})\longmapsto
 (a_{\rm free},-\textstyle\sum_l a_{jl}\bar g_l(z))_{j=1}^s,
\] the free-coordinate derivative is the identity. Choose \(s\)
independent coordinate differential combinations for the remaining rows;
then the remaining block has rank \(s\). Some full minor is nonzero, so
the output differentials are independent in \(k(Z_i)(A_{\rm free})\). In
characteristic zero, the rank of differentials of a finitely generated
subfield equals its transcendence degree (choose a separating
transcendence basis). Thus the output functions are algebraically
independent. Their embedding into
\(\widetilde{\mathcal H(y_i)}(A_{\rm free})\) preserves this
independence.

Every coefficient polynomial over \(C\) now has its exact full Gauss
norm after evaluation: normalize by a coefficient of maximum norm and
reduce; the resulting nonzero polynomial cannot vanish at the
independent residues. Fractions and completion give isometric embeddings
of the \textbf{same} field \[
 F=\widehat{C(A_{j0},\ldots,A_{jN}:1\leq j\leq s)}
 \hookrightarrow H'_i. \tag{6.1}
\] Completion changes neither the value group nor residue field:
approximate each nonzero element with error strictly below its norm,
then use the same norm and residue. The field \(F\) is a transcendental
Gauss extension, not a subfield of \(\widehat{\overline C}\).

Over the honest rational coefficient field, use the independent generic
hyperplanes \[
 A_{j0}S_0+\sum_{l=1}^N A_{jl}S_l=0\quad(1\leq j\leq s).
 \tag{6.2}
\] They cut an actual geometrically smooth curve. For any smooth
embedded geometric component of dimension \(e>0\), vanishing at a point
and vanishing differential impose \(e+1\) independent linear hyperplane
conditions: very ampleness makes the affine coordinate ratios span
cotangents. The bad incidence in point times projective hyperplane space
has dimension at most \(N-1\). Projectivity makes its hyperplane image
closed in the \(N\)-dimensional parameter space, so the generic
hyperplane avoids it, including after algebraic closure. A generic
hyperplane contains no positive-dimensional component, and ampleness
makes its intersection with each such component nonempty. It cuts
dimension by one. Apply this successively to the independent rows; each
step is smooth and the equations form a regular sequence.

For geometric integrality, the Koszul resolution has terms
\(\mathcal O(-kM)^{\binom{s}{k}}\) in degree \(-k\). In its
hypercohomology sequence, the terms of total degree zero are
\(H^k(\mathcal O(-kM))\), and potential incoming terms have
\(H^{k-1}(\mathcal O(-kM))\). All vanish for \(k>0\) by (4.1), together
with \(H^0(\mathcal O(-kM))=0\) for negative ample powers. For the
latter, a nonzero section of \(\mathcal O(-kM)\) would inject
\(H^0(\mathcal O(kM))\) into \(H^0(\mathcal O_X)=C\) on the integral
variety. But \(kM\) is very ample on a positive-dimensional variety, so
its section space has dimension at least two. Proper geometric
integrality gives \(H^0(\mathcal O_X)=C\): over an algebraic closure a
regular function has closed image in \(\mathbb P^1\) missing infinity
and hence is constant, and flat field Cech base change descends that
statement. There is no outgoing differential from the degree-zero term
to positive complex degree, which is absent. Thus
\(H^0(\mathcal O_D)=F\). To justify this after any field extension, use
a finite affine cover of the separated projective scheme: intersections
are affine and its quasi-coherent Cech complex tensors with the flat
field extension. Its cohomology therefore commutes with that extension.
The same Koszul computation holds over an algebraic closure and then the
completion. A smooth geometric projective curve has disjoint irreducible
components; more than one gives a nontrivial global idempotent. Hence it
is geometrically integral.

We obtain an actual smooth geometrically integral projective curve
\(D/F\) in the actual \(X_F\). With \(A=L|_D\), its raw degree is \[
 a=\deg A=(L\cdot M^{n-1})=d^{n-1}V. \tag{6.3}
\]

The coordinate evaluations define honest analytic points of \(D\) over
\(F\), since the equations (6.2) vanish. Their residue contains the
original coordinate subfield and the free coefficient residues, of
transcendence degree \(n+sN\) over \(k\). The base residue has
transcendence degree \(s(N+1)\), so relative transcendence degree is at
least one. A nongeneric point of the integral curve is closed, its
scheme field is finite over \(F\), and its valued residue is algebraic
over \(\widetilde F\): independent residue lifts would be linearly
independent over \(F\), by maximal-coefficient reduction. Completion
changes no residues. Thus the analytic points have \textbf{generic curve
support before} invoking any embedding of \(F(D)\).

Now \(F(D)\) embeds honestly in \(H'_i\). It contains the evaluations of
\(C(X)\) and all \(A_{\rm free}\). The first is dense in
\(\mathcal H(y_i)\); approximation of finitely many
numerator/denominator coefficients, with denominator error below its
nonzero norm, makes \(F(D)\) dense in \(H'_i\). Consequently the
completed residue field at this actual point is \(H'_i\), and its
relative residue transcendence degree is exactly one.

For the value group, the elementary Abhyankar independence argument says
that \(d\) independent unit residues and \(q\) rationally independent
values modulo the base group yield \(d+q\) algebraically independent
elements. Group a polynomial by the value-exponent vectors; independent
residues prevent cancellation within each group, and distinct value
classes prevent cancellation between groups. Apply this to the valued
\(C(X)\): \(d=n\) forces rational rank zero. Because the group of \(C\)
is divisible, a real ordered torsion extension has no new values. Unit
Gauss variables and completion retain this same group.

Set \(B=\widehat{\overline F}\). An algebraic-closure embedding over the
complete \(F\) into an algebraic closure of \(H'_i\) is isometric by
uniqueness of valuation extension; it extends to
\(B\hookrightarrow\Omega_i=\widehat{\overline{H'_i}}\). Both residue
extensions are algebraic, and their value groups are the same divisible
group. The original coordinate evaluations define actual
\(z_i\in D_B^{\rm an}\). The residual element transcendental over
\(\widetilde F\) can be lifted from \(F(D)\) by the density just proved;
it remains transcendental after the algebraic residue extension to
\(\widetilde B\). Hence \(z_i\) has generic support before invoking
\(B(D)\). Its residual transcendence degree is one and its relative
value rank zero, so it is type II. The completed residue field of
\(z_i\) need not equal \(\Omega_i\); that equality is neither asserted
nor needed. Every original rational-function value is preserved in these
embeddings. An original section norm in a pulled-back original reference
is consequently preserved too.

For \(n=1\), take \(B=C\), \(D_B=X_C\), and \(z_i=y_i\) directly. Their
maximal residual degree one and unchanged divisible group make them type
II. No incidence construction or inverse curve equation is needed. GM
A.7 is used nowhere for the transcendental extension (6.1); its
algebraic-completion restriction is respected.

\subsection{Curve section degrees bound all slopes on fixed
paths}\label{curve-section-degrees-bound-all-slopes-on-fixed-paths}

Choose once a semistable vertex set of the actual smooth projective
\(B\)-curve, include every \(z_i\) by BPR3.13, and let \(\Gamma\) be its
compact connected skeleton. Keep every cycle. Put \[
 D_0=\max_{i,j}d_\Gamma(z_i,z_j)<\infty,
\] with \(D_0=0\) for one point. These data are frozen before any later
FS family.

\textbf{Lemma 3 (finite extension after completion).} If a valued
extension \(E/E_0\) has finite degree \(\delta\), then its completion
\(H\) has degree at most \(\delta\) over the completed base \(F_0\). Its
residue degree is at most \(\delta\).

\textbf{Proof.} Embed \(F_0\) isometrically in \(H\). The \(F_0\)-span
of an \(E_0\)-basis has dimension at most \(\delta\) and is closed. Here
completeness of finite-dimensional normed spaces requires no spherical
completeness. Induct on an \(F_0\)-basis: the preceding span \(W'\) is
complete and closed, and for a new vector \(v\notin W'\),
\(c=\inf_{w\in W'}|v-w|>0\). For \(b\ne0\), \[
 |w+bv|\geq |b|c,\qquad
 |w|\leq\max(|w+bv|,|b||v|).
\] These give equivalence with the maximum coordinate norm and
completeness. A closed span containing dense \(E\) equals \(H\).
Independent residue unit lifts are \(F_0\)-linearly independent: divide
any finite relation by a coefficient of maximal norm and reduce. Thus
residue degree is bounded by vector-space degree. No defectlessness or
equality in the fundamental inequality is assumed. \(\square\)

Take any nonzero \(t\in H^0(D_B,kM_D)\) and \(f=t/S_0^k\). For
nonconstant \(f\), its finite map to \(\mathbb P^1\) has degree \[
 \delta=[B(D):B(f)]=\deg(f)_\infty
       \leq k\deg(M_D)=kd\,a. \tag{7.1}
\] The pole divisor is coefficientwise bounded by
\(k\operatorname{div}(S_0)\), and cancellations only decrease it. The
degree identity follows by pulling back infinity; the finite flat curve
map and the effective Cartier degree formula apply to the actual curve.
Constant \(f\) has zero slope.

Refine the skeleton after deleting the finite zeros and poles of \(f\),
retaining the old \(\Gamma\) with its original edge lengths. On each
resulting annular edge portion, \(f\) is an analytic unit. On a closed
strict subannulus the unit form is \[
 f=cT^q(1+g),\qquad q\in\mathbb Z,\quad\|g\|_{\sup}<1. \tag{7.2}
\] BPR2.2(1) states it. An immediate independent explanation is useful:
at a value-group radius, the normalized Laurent reductions of a unit and
its inverse multiply to one in \(\widetilde B[\tau,\tau^{-1}]\), so each
is a monomial. Thus exactly one Laurent exponent dominates there. A
change of dominant exponent would give a tie at a radius in the
divisible value group, contradicting the monomial reduction. Convergence
of both Laurent series leaves only finitely many potentially dominant
exponents on the compact interval. It gives a uniform strict gap for the
remaining terms, proving (7.2). This does not use the inverse curve MA
equation.

Every point of the compact curve skeleton has generic scheme support:
over algebraically closed \(B\), a nongeneric curve support is a
\(B\)-rational closed point of type I, while the skeleton contains only
types II and III. At a dense value-group Gauss radius \(\rho=|b|\),
write \(\tau=T/b\). Its actual completed residue field has residue
\(\widetilde B(\tau)\) and generic algebraic curve support. If
\(q\ne0\), the residue of \(f/(cb^q)\) is \(\tau^q\).
Maximal-coefficient reduction proves that every polynomial in this
normalized \(f\) has its exact Gauss norm, so the completed subfield of
honest \(B(f)\) has residue \(\widetilde B(\tau^q)\). Its index in
\(\widetilde B(\tau)\) is \(|q|\). Indeed Laurent exponents have unique
residue classes modulo \(|q|\), giving the usual rational-function basis
and independence. Lemma3 applied to the finite honest extension
\(B(D)/B(f)\) now gives \[
 |q|\leq\delta\leq kd\,a. \tag{7.3}
\] The case \(q=0\) has zero slope. Density of the divisible value-group
radii and continuity give the estimate at type-III limits and closed
endpoints. Integrating along any old shortest path of \(\Gamma\) yields
\[
 |\log|f(z_i)|-\log|f(z_j)||
       \leq kd\,a\,d_\Gamma(z_i,z_j). \tag{7.4}
\] BPR3.13 subdivisions and attached tails retain that path and its
length. No cycle is removed and no varying graph diameter is used.

\subsection{Every FS family and every actual
endpoint}\label{every-fs-family-and-every-actual-endpoint}

Let \(\psi\) be any original FS metric on the actual \(L\), with all
real weights permitted. By taking powers of its finite nonzero defining
sections, align its denominator with \(d\); multiply the weights by the
same powers. Every original nonzero section restricts nonzero to
\(D_B\), since its nonzero generic value at each original \(y_i\) is
preserved at \(z_i\). In the rational \(S_0/d\) frame its weight is \[
 w_\psi=\frac1{kd}\max_l\bigl(\log|t_l/S_0^k|+\beta_l\bigr). \tag{8.1}
\] The entire expression, including every weight, is divided by \(kd\).
Each term is \(a\)-Lipschitz on the fixed skeleton by (7.4), and a
finite maximum preserves that constant. At every \(z_i\), the
reference's frame weight is zero: \(S_0\) occurs, its ratio is one, and
all other original ratios have norm at most one. These are the same
values as at \(y_i\).

Normalize \(\sup_T(\psi-\theta)=0\). D's surjective completed-closure
projection gives the same supremum upstairs. Equation (4.3) norms every
defining section; commuting finite maxima and real weights shows some
\(y_j\) has scalar zero. Thus \[
 -aD_0\leq(\psi_C-\theta_C)(y_i)\leq0
 \quad\text{for every }i\text{ and every original FS metric}. \tag{8.2}
\] All initial geometric data precede \(\psi\).

Now take \textbf{any actual nontrivial psh} scalar \(u\), before
assigning finite-energy membership, and a defining FS net
\(u_\lambda\downarrow u\). Put \(s_\lambda=\sup_Tu_\lambda\),
\(s=\sup_Tu\). These decrease to \(s\): otherwise directed nonempty
closed sets \(\{u_\lambda\geq s+\epsilon\}\) would have nonempty
intersection by compactness, contradicting \(\sup u=s\). Nontriviality
and boundedness above make \(s\) finite. Apply (8.2) to
\(u_\lambda-s_\lambda\) and pass at each of the finitely many points to
obtain \[
 -aD_0\leq u_C(y_i)-s\leq0. \tag{8.3}
\] In particular these endpoint values are finite. This proof uses
neither weighted E1 continuity nor an assumed finite-energy endpoint.

For \(s=0\), the genuine atomic reference measure (5.1) gives \[
 \int(-u_C)\,\operatorname{MA}_{\theta_C}(0)
       \leq VaD_0.
\] O's raw continuous-field pushforward and the Borel pullback identity
identify this with the original reference integral. This proves (2.1)
for the fixed \(\theta\), with \(M_0=VaD_0\). Neither inverse atoms nor
individual raw coefficients have entered the proof. At trivial valuation
T's canonical norming generic point supplies moment zero and its
complete theorem supplies the endpoint equation.

\subsection{All-test integrability, a unit bound, and finite
dual}\label{all-test-integrability-a-unit-bound-and-finite-dual}

This section needs only O and the established moment bound for the
implication \(U<\infty\Rightarrow J<\infty\).

\textbf{Lemma 4 (uniform unit bound without an anchor).} Every actual E1
potential is absolutely \(\mu\)-integrable if and only if
\(U_\mu<\infty\).

\textbf{Proof.} If all tests are integrable but \(U=\infty\), choose
normalized \(u_j\) of cost at most one with \(A_\mu(u_j)\geq4^j\). The
actual convex sums \[
 w_N=\sum_{j=1}^N2^{-j}u_j+2^{-N}0
\] decrease, are nonpositive, and satisfy \(E(w_N)\geq-1\) by concavity.
Their all-point endpoint is in the decreasing closure. It cannot be
identically \(-\infty\): for any \(R>0\), the directed closed sets
\(\{w_N\geq-R\}\) would have empty intersection, hence a tail would lie
below \(-R\), and energy monotonicity would force \(E(w_N)\leq-RV\), a
contradiction. Use the extended majorant energy at this stage, before E1
membership. The nontrivial endpoint \(w\) is actual psh, and O's
decreasing energy gives \(E(w)\geq-1\), so it is E1. Ordinary
\textbf{sequential} monotone convergence gives \[
 \int(-w)\,d\mu=\sum_j2^{-j}A_\mu(u_j)\geq\sum_j2^j=\infty,
\] a contradiction. Conversely energy concavity gives
\(B(tu)\leq tB(u)\). For normalized \(u\), choose \(t=1\) if \(B\leq1\),
and \(t=1/B\) otherwise. Then \(t>0\), \(B(tu)\leq1\), and
\(tA_\mu(u)\leq U<\infty\). Finite constant shifts give absolute
integrability for every unnormalized E1 scalar. \(\square\)

\textbf{Lemma 5 (the exact Radon-net law and finite continuous dual).}
If \(J_\theta(\mu)<\infty\), every actual E1 test is integrable. For a
continuous psh net \(v_\lambda\downarrow u\in\mathcal E^1\), both its
energy and target integral converge to those of \(u\). The continuous
and all-E1 dual suprema agree.

\textbf{Proof.} On a tail take a finite upper bound \(R\) for all
\(v_\lambda\). The dual inequality gives \[
 \int(R-v_\lambda)\,d\mu
   \leq RV+J-E(v_\lambda)\leq RV+J-E(u)<\infty.
\] For a directed increasing net of nonnegative continuous, or lower
semicontinuous, functions on a compact Hausdorff space with a Radon
measure, its integral equals the supremum of stage integrals. To prove
the precise law, each strict open superlevel of the limit is the
directed union of stage superlevels. A compact inner approximation is
covered by finitely many of these opens, hence one common index covers
it. This identifies the mass at each level as the supremum of stage
masses. For finitely many levels choose one common index approximating
all masses; finite lower layer sums then converge. Truncate height and
refine layers to recover the full integral, even if infinite. Arbitrary
measurable-net monotone convergence is not assumed.

Apply this law to \(R-v_\lambda\). The uniform finite bound excludes an
infinite endpoint integral. It also proves integral convergence, while O
gives \(E(v_\lambda)\downarrow E(u)\). Hence
\(E(u)-\int u\,d\mu\leq J\). Continuous tests are E1, so the suprema
agree. \(\square\)

\textbf{Lemma 6 (Bernstein scaling with the correct constant).} For
normalized actual \(u\in\mathcal E^1\) and every real \(0\leq t\leq1\),
\[
 B(tu)\leq tM_0+D_nt^2B(u),\qquad D_n=\frac{n(n+1)}2. \tag{9.1}
\]

\textbf{Proof.} O makes \[
 A_k=\int(-u)\,\operatorname{MA}(u^{[k]},0^{[n-k]})
\] finite and nonnegative. Its genuine energy identity gives
\(\sum_kA_k=(n+1)B(u)\), and \(A_0\leq M_0\). The affine curvature law
and actual segment derivative yield the finite polynomial identity \[
 B(tu)=\sum_{k=0}^n\binom nk A_k
          \int_0^t s^k(1-s)^{n-k}\,ds. \tag{9.2}
\] For \(k=0\) the integral is at most \(t\). For \(k\geq1\), the
nonnegative Bernstein weight \(\binom nk s^k(1-s)^{n-k}\) is at most the
sum of all positive-index weights, \(1-(1-s)^n\leq ns\). Hence their
contribution is at most \((n+1)B(u)\int_0^t ns\,ds=D_nt^2B(u)\). No
integral is undefined; zero coefficients at poles are omitted in the
reference endpoint. \(\square\)

Assume now \textbf{\(U_\mu<\infty\)}. Choose \[
 t=(1+2M_0+\sqrt{2D_nB(u)})^{-1}>0.
\] The first term \(tM_0\) in (9.1) is at most \(1/2\), and the second
at most \(1/2\); if \(B=0\) the latter is zero. Thus \(B(tu)\leq1\). Its
supremum is zero, so \(tA_\mu(u)\leq U\). This proves (2.2) and then
(2.3). Completing the square in \(\sqrt B\) gives (2.4). Lemma4 and
Lemma5 complete the equivalence. Every coercivity conclusion in this new
manuscript is quantified by the finite target conditions.

For two continuous psh references, O's cocycle gives a constant energy
shift and a bounded scalar change. Consequently the absolute E1 class,
all-test integrability, and finite continuous dual are
reference-independent. The selected \(M_0\) proof is not extended by
pretending the new reference has the old support.

\subsection{What this proves toward attainment, and what
remains}\label{what-this-proves-toward-attainment-and-what-remains}

\textbf{Proposition 7 (an already known output attains its own dual).}
If \(\mu=\operatorname{MA}(u)\) for actual \(u\in\mathcal E^1\), then
its intrinsic criterion holds and \(u\) maximizes its actual
energy-minus-integral functional.

\textbf{Proof.} O gives mass \(V\) and all-test integrability, so
Theorem2 gives finite \(J\). Concavity and the finite segment derivative
at \(u\) give \[
 E(v)-E(u)\leq\int(v-u)\,\operatorname{MA}(u)
\] for every actual E1 \(v\). All terms are integrable; differences mean
their common finite-value full-measure representatives, without
subtracting poles. Therefore \(F_\mu(v)\leq F_\mu(u)\), and Lemma5
recovers equality with the continuous-test dual. This is a supporting
inequality for a known genuine output; it supplies no inverse solution.
\(\square\)

For any intrinsic target, coercivity forces a normalized maximizing net
into a bounded energy-cost region. A quantitative bound is available.
Put \(C=U(1+2M_0)\) and \(c=U\sqrt{2D_n}\). On a tail with
\(F_\mu(u)\geq-1\), (2.3) implies \[
 B(u)\leq
 \left(\frac{c+\sqrt{c^2+4(C+1)}}2\right)^2. \tag{10.1}
\] This identifies the precise region in which an actual compactness or
attainment replacement is required. It proves neither compactness of
that region nor upper semicontinuity of \(F_\mu\) in an available
topology.

A possible continuation is to take product-topology limits of normalized
potentials or their section norms. The bound \([-\infty,0]\) supplies
abstract coordinate limits, and (8.3) excludes the identically
\(-\infty\) coordinate vector. It does not make a limit upper
semicontinuous, actual psh, or an all-point decreasing-FS endpoint.
Abstract energy completion has the same realization obligation. Weak
compactness of Radon measures controls measures, not their actual
potentials. None is used as a replacement theorem.

O proves derivatives along actual E1 segments. It does not prove
derivatives of \(E(P(u+tf))\) for an arbitrary continuous outer
direction \(f\), for both signs of \(t\), nor envelope orthogonality.
BE7.24--30 distinguishes continuous obstacles, psh envelopes, and
continuity of envelopes. The general-field envelope step is not obtained
by deleting its hypothesis or by invoking the continuous
scalar-extension theorem. C's comparison concerns strict inequalities;
it does not identify every solution at every point, including common
poles. Equality of finite-energy norms or of MA measures also does not
by itself recover the all-point scalar.

The original equation obligations therefore remain:

\begin{enumerate}
\def\labelenumi{\arabic{enumi}.}
\tightlist
\item
  Prove normalized compactness of actual psh potentials in a topology
  with the necessary energy and target-integral semicontinuity, or prove
  an adequate actual realization/attainment substitute on (10.1).
\item
  Prove all-point envelopes, orthogonality, and both-sign outer-obstacle
  differentiation where the chosen variational argument uses them, or
  replace each use with a complete argument.
\item
  Attain every prescribed positive Radon target of raw mass \(V\)
  satisfying the intrinsic criterion.
\item
  Show that every such maximizer solves the genuine equation and prove
  equality of every pair of actual E1 solutions at every analytic point
  modulo one finite real constant.
\end{enumerate}

T and whole G227 already discharge these obligations in their exact
trivial-field and original-field genus-zero domains. They do not
discharge them for positive-genus curves over nontrivial fields or in
general dimension. The alternative positive-genus formal realization and
kernel route still has unresolved comparison/source/model antecedents,
including review198 premises14/15. No recognition theorem with those
unestablished premises is asserted here. A targeted public literature
discovery search also found Reboulet's arbitrary-field geodesic theorem
with an explicit continuity-of-envelopes assumption
(\url{https://arxiv.org/abs/2012.07972}) and a degeneration result
restricted to \(\mathbb C((t))\)
(\url{https://jep.centre-mersenne.org/articles/10.5802/jep.229/}). Their
abstracts are scope checks only, not new mathematical premises. The
search is not exhaustive and no absence-of-literature theorem is
claimed. The remaining original P16.4 is open.

\subsection{Evidence, reruns, and
limitations}\label{evidence-reruns-and-limitations}

\texttt{prepare.py} binds all seven complete grants and unchanged
clarifications, preserves the failed data units, and checks the 26
frozen primary files against the supplied byte lengths and SHA256
values. The original \texttt{input.json} is unchanged.
\texttt{read\_source.py\ LABEL\ FIRST\ LAST} uses literal LF numbering
and appends exact source-read bindings to \texttt{source-reads.jsonl}.
Source reading is evidence of what was inspected, not mathematical
acceptance. Public immediate Stacks response records are retained in
\texttt{data/web1.json} through \texttt{data/web3.json}; the application
ledger is quote-free. No private source bytes are published by this
work.

\texttt{finite\_checks.py} uses exact rational arithmetic. It verifies
492 dimension/time Bernstein pairs (\(n=1,\ldots,12\), \(t=j/40\)), 324
endpoint coefficients, 240 scaling cases, 60 raw-mass counterexample
formula rows, 2,904 Laurent exponent decompositions, a 625-case two-row
rank grid (496 nonsingular,129 singular), 552 Koszul diagonal terms,
1,296 fixed-cycle path checks, and 11,664 weighted-maximum path checks
with both weights in \(\{-5/2,0,7/3\}\) and the entire maximum divided
by seven. Exhaustiveness applies only to those domains. The source-based
geometric existence, true MA identification, arbitrary real weights, and
all-field analytic claims are mathematical arguments, not finite
exhaustive results. The counterexample rows instantiate formulas whose
genuine metric interpretation is proved in Proposition1; arithmetic
alone does not establish that interpretation.

Rerun in this directory: \texttt{python3\ prepare.py};
\texttt{python3\ finite\_checks.py};
\texttt{python3\ build\_deliverable.py};
\texttt{python3\ emit\_notes.py}. The standalone source
\texttt{P16.4-paper.tex} is generated from this full argument and
compiled with existing pandoc/pdfLaTeX, with shell escape disabled.
\texttt{commands.json}, \texttt{environment.json}, build logs,
\texttt{build-result.json}, and \texttt{artifact-hashes.json} retain
rerun evidence. There was no exposed Xean native execution service or
supervised compute service; no model-specific remote invocation is
claimed. No resource-heavy calculation or installation was launched. The
new results are unverified drafts and \texttt{candidate=false} is
required.

\subsection{Primary references}\label{primary-references}

BE: S. Boucksom and D. Eriksson, \emph{Spaces of norms, determinant of
cohomology and Fekete points in non-Archimedean geometry}, literal
arXiv:1805.01016v3, \url{https://arxiv.org/pdf/1805.01016v3}.

CLD: A. Chambert-Loir and A. Ducros, \emph{Formes differentielles
reelles et courants sur les espaces de Berkovich}, literal
arXiv:1204.6277v1, \url{https://arxiv.org/pdf/1204.6277v1}.

GRW: W. Gubler, J. Rabinoff and A. Werner, \emph{Tropical skeletons},
original publisher version,
\url{https://aif.centre-mersenne.org/item/10.5802/aif.3125.pdf}.

Poineau: J. Poineau, \emph{Les espaces de Berkovich sont angelique},
original publisher version,
\url{https://www.numdam.org/article/BSMF_2013__141_2_267_0.pdf}.

GM: W. Gubler and F. Martin, \emph{On Zhang's semipositive metrics},
Documenta Mathematica24 (2019),331--372,
\url{https://ems.press/content/serial-article-files/26466}.

BPR: M. Baker, S. Payne and J. Rabinoff, \emph{On the structure of
nonarchimedean analytic curves}, literal arXiv:1404.0279v1,
\url{https://arxiv.org/pdf/1404.0279v1}.

FK2017: K. Fujiwara and F. Kato, \emph{Foundations of Rigid Geometry I},
literal arXiv:1308.4734v5, \url{https://arxiv.org/pdf/1308.4734v5}. Only
the ordinary restricted-series closed-ideal unit0.8.4.4 is used for the
formal-presentation application. It is distinct from Freitag--Kiehl and
from the unused formal GAGA/comparison route.

Stacks Project: exact units and immediate proof applications are listed
in the source ledger, with tags0892,0B5T,0FVV,0BXX,0AYY,0AYZ. Complete
caller-grant bodies and all unchanged clarifications are bound to the
assignment in \texttt{data-bindings.json}; the seven supplied grant
labels have the scope stated at the start.

\end{document}
